\section{AI辅助流程及评价指标}
\subsection{受约束交互与证据链}
流程如图\ref{fig:paper-workflow}所示。AI辅助编写和审查参数生成、结果提取与证据核验程序；自然语言入口接受约定中文模板，映射到粒子、方向、尺寸和存储位等结构化参数，再检查范围与单位。R50演示支持预览、审计和定位已有证据，实际求解由冻结输入和固定控制器执行。该结果证明受约束请求的参数映射，不能写成任意自然语言生成任意模型或自主开启新求解。

\begin{figure}[htbp]
\centering
% Conceptual workflow, not a measured result or screenshot.
\resizebox{\linewidth}{!}{%
\begin{tikzpicture}[
 block/.style={draw=signalblue,rounded corners=2pt,fill=softgray,
 text width=4.2cm,minimum height=1.5cm,align=center,font=\small},
 arrow/.style={-{Latex[length=2.4mm]},thick,draw=signalblue},
 node distance=0.7cm and 0.6cm]
\node[block] (request) {限定中文请求\\尺寸、粒子、角度\\独立事件数};
\node[block,right=of request] (matrix) {受约束的工况清单\\单位、范围、版本与种子};
\node[block,right=of matrix] (transport) {Geant4 硅内电离沉积\\$E_{\rm ion}\ \rightarrow\ Q_{\rm dep}$\\有限径迹及单位检查};
\node[block,below=of transport] (tcad) {TCAD 器件与电路求解\\等效径迹与受控初态\\读出瞬态与端口检查};
\node[block,left=of tcad] (metric) {匹配中性参考与同位对照\\$D,\ M,\ \Delta Q,\ \Delta M$};
\node[block,left=of metric] (model) {按工况分组的响应拟合\\冻结预测后检验未见响应};
\draw[arrow] (request)--(matrix);
\draw[arrow] (matrix)--(transport);
\draw[arrow] (transport)--(tcad);
\draw[arrow] (tcad)--(metric);
\draw[arrow] (metric)--(model);
\node[below=0.40cm of metric,align=center,font=\small,text width=14.5cm]
 {每一层绑定输入版本、运行记录和来源哈希；科学质量门槛与执行状态分别记录。};
\end{tikzpicture}}

\caption{AI辅助评估流程及证据传递。自然语言映射、粒子事件、TCAD响应和冻结预测具有不同产物；本文只将已归档并通过相应检查的数据作为结果。}
\label{fig:paper-workflow}
\end{figure}

每个响应关联模型版本、粒子事件和输入标识。检查覆盖初态、预充、源项积分、端口电流闭合以及所需匹配关系；求解结束标志不足以证明物理评价可用。确定性重跑保留原失败身份并关联替代归档，不增加独立样本。流程将AI建议、固定规则与求解器结果分开，便于定位错误来源；本文未设人工操作对照组，因此不报告AI带来的效率提升比例。

\subsection{粒子增量与读出参考的分工}
分别用$e$、$c$、$n$表示粒子事件、同条件同存储位无粒子对照和中性参考。在保持阶段$t_h=2.8\ns$，定义
\begin{equation}
\Delta Q=C_s\big[V_{{\rm SN},e}(t_h)-V_{{\rm SN},c}(t_h)\big].
\label{eq:paper-dq}
\end{equation}
$\Delta Q$以fC报告，负值表示存储电容相对对照减少电荷。两种位值对同一方向扰动的敏感性不同，因此不能把全部$\Delta Q$称为正的电荷损失。

对位值$b\in\{0,1\}$，在$t_r=4\ns$定义
\begin{align}
D_e&=V_{{\rm BL},e}(t_r)-V_{{\rm BL},n}(t_r),\\
M_e&=(2b-1)D_e,\qquad
M_c=(2b-1)\big[V_{{\rm BL},c}(t_r)-V_{{\rm BL},n}(t_r)\big],\\
\Delta M&=M_c-M_e.
\label{eq:paper-margin}
\end{align}
$D$、$M$、$\Delta M$以mV报告，$\Delta M>0$表示粒子引起裕量下降。中性参考定义读出差分，位匹配对照提取粒子增量；二者不可互换。同一匹配参考在$\Delta M$中相消，也不意味着可以省略它，因为绝对裕量和判据分类仍取决于参考。

主判据阈值为$\epsilon=5\mV$：$M>\epsilon$为正确方向，$M<-\epsilon$为判据误读，$|M|\leq\epsilon$为临界；另按$2.5/10\mV$检查阈值敏感性。阈值为项目约定，未标定实际感放失调和噪声。数值检查不足以支持边界分类的响应保留“数值不确定”，不能并入正确或误读。本文的误读是固定采样时刻的模型判据，不等同于已锁存的持久位翻转。

\subsection{矩阵设计与独立事件单位}
设计包含4种栅长、2种粒子、2个角度，每工况4个独立事件，共64事件。每个事件分别计算存储0和1，共128个配对响应。训练采用$L=0.18/0.35\um$，含32事件、64响应；$L=0.25\um$为尺寸插值留出，$L=0.45\um$为尺寸外推，各16事件、32响应。训练与测试按尺寸分离，配对位值及同事件的氧化层或数值变体均不增加粒子事件分母。该设计检验特定尺寸变化，不构成对任意能量、入射位置或材料体系的泛化验证。
